% ID: SEL-D0049-OA-E04
% ESTADO: validada. S050.
\textbf{a)} La proporción muestral es $\hat p=240/400=0{,}60$. Con $240$ éxitos y $160$ fracasos usamos la aproximación normal. Para un $98{,}5\%$ de confianza queda $0{,}0075$ en cada cola y $z_{0{,}9925}\approx2{,}4324$. Así,
\[
0{,}60\pm2{,}4324\sqrt{\frac{0{,}60\cdot0{,}40}{400}}
\approx\boxed{[0{,}5404,\ 0{,}6596]}.
\]
\textbf{b)} La amplitud es $2z_{0{,}9925}\sqrt{\hat p(1-\hat p)}/\sqrt n$. Con el numerador fijo, aumenta al reducir $n$ y disminuye al aumentar $n$. Por ejemplo, cuadruplicar $n$ reduce la amplitud a la mitad.
